\chapter{The stochastic Bornhuetter-Ferguson model}

Mack, \emph{The Prediction Error of Bornhuetter/Ferguson} (2008), states the
increment model on page 91 and derives the prediction-error calculation on
pages 95--99. The exact probability statements and the practical assessments
are separated below.

\begin{definition}[BF1--BF3 and the future reserve]\label{def:bfModel}\lean{VerifiedReserving.RandomTriangle.bfIncrementRv, VerifiedReserving.BFFullIncrementIndependence, VerifiedReserving.BFIncrementMean, VerifiedReserving.BFIncrementVariance, VerifiedReserving.BFPatternNormalized, VerifiedReserving.bfCumulativePattern, VerifiedReserving.RandomTriangle.bfFutureReserveRv}\leanok
For zero-based increments $S_{ik}$, BF1 says that the full family is
independent. On a finite horizon $m$, BF2 and BF3 say
\[
 E[S_{ik}]=x_i y_k,\qquad \operatorname{Var}(S_{ik})=x_i\sigma_k^2,
 \qquad \sum_{k<m}y_k=1.
\]
The true reserve after the first $d$ periods is
$R_i=\sum_{d\leq k<m}S_{ik}$, and
$z_d=\sum_{k<d}y_k$ is the expected reported proportion.
\end{definition}

\begin{theorem}[Mean and process variance of a BF reserve]\label{thm:bfModelMoments}\lean{VerifiedReserving.integral_bfFutureReserveRv_eq_tail, VerifiedReserving.integral_bfFutureReserveRv, VerifiedReserving.variance_bfFutureReserveRv}\leanok
\uses{def:bfModel}
Under BF2 and integrability,
\[
 E[R_i]=x_i\sum_{d\leq k<m}y_k=x_i(1-z_d).
\]
Under BF1, BF3, and square integrability,
\[
 \operatorname{Var}(R_i)=x_i\sum_{d\leq k<m}\sigma_k^2.
\]
\end{theorem}
\begin{proof}\leanok The integral commutes with the finite sum. Pairwise
independence supplied by BF1 makes the finite-sum variance additive. \end{proof}

\begin{definition}[Random BF estimate and MSEP]\label{def:bfRandomEstimate}\lean{VerifiedReserving.bfStochasticReserveEstimate, VerifiedReserving.bfMeanSquaredPredictionError, VerifiedReserving.BFObservedPredictionIndependence}\leanok
For a random initial ultimate $U$ and a random selected cumulative pattern $z$,
the BF estimate is $\widehat R^{BF}=U(1-z)$. The unconditional MSEP is
$E[(\widehat R^{BF}-R)^2]$. Mack's source conditions on the observed increments;
common independence means that the pair consisting of estimate and future
reserve is independent of the observed-data sigma-algebra.
\end{definition}

\begin{theorem}[Product variance and the single-row MSEP]\label{thm:bfPredictionError}\lean{VerifiedReserving.variance_mul_of_indepFun, VerifiedReserving.variance_bfStochasticReserveEstimate, VerifiedReserving.bfMeanSquaredPredictionError_eq_variance_add, VerifiedReserving.bfMeanSquaredPredictionError_eq, VerifiedReserving.condExp_sq_bfPredictionError_eq}\leanok
\uses{def:bfModel, thm:bfModelMoments, def:bfRandomEstimate}
If $U$ and $z$ are independent and the displayed products are
square-integrable, then
\[
 \operatorname{Var}\bigl(U(1-z)\bigr)
 =\bigl(E[U]^2+\operatorname{Var}(U)\bigr)\operatorname{Var}(z)
  +\operatorname{Var}(U)(1-E[z])^2.
\]
If the BF estimate and true future reserve are independent and have the same
mean, then
\[
 \operatorname{msep}(\widehat R^{BF})
 =\operatorname{Var}(\widehat R^{BF})
  +x_i\sum_{d\leq k<m}\sigma_k^2.
\]
Under common independence from an observed-data sigma-algebra, the conditional
expectation of the squared prediction error equals that same constant almost
surely.
\end{theorem}
\begin{proof}\leanok Factor first and second product moments by independence.
For MSEP, center the prediction error and add the variances of the independent
estimate and target. \end{proof}

\begin{definition}[Raw pattern and variance estimators]\label{def:bfRaw}\lean{VerifiedReserving.RandomTriangle.bfYRawRv, VerifiedReserving.RandomTriangle.bfLinearPatternRv, VerifiedReserving.RandomTriangle.bfResidualSumSqRv, VerifiedReserving.RandomTriangle.bfSigma2RawOnRv}\leanok
On a finite set $A$ of accident years with known volumes $x_i$, equations~(1)
and~(2) become
\[
 \widehat y_k=\frac{\sum_{i\in A}S_{ik}}{\sum_{i\in A}x_i},\qquad
 \widehat\sigma_k^2=\frac1{|A|-1}\sum_{i\in A}
   \frac{(S_{ik}-x_i\widehat y_k)^2}{x_i}.
\]
A general linear estimator is $\sum_{i\in A}a_iS_{ik}$.
\end{definition}

\begin{theorem}[Mack equations (1) and (2)]\label{thm:bfRaw}\lean{VerifiedReserving.integral_bfYRawRv, VerifiedReserving.variance_bfYRawRv, VerifiedReserving.integral_bfLinearPatternRv, VerifiedReserving.variance_bfLinearPatternRv, VerifiedReserving.bfLinearPatternRisk_sub_raw, VerifiedReserving.bfYRawRv_best_linear_unbiased, VerifiedReserving.bf_raw_weighted_sq_dev, VerifiedReserving.RandomTriangle.bfResidualSumSqRv_eq, VerifiedReserving.integral_bfSigma2RawOnRv}\leanok
\uses{def:bfModel, def:bfRaw}
For nonnegative volumes with nonzero total, the raw estimator is unbiased and
has variance $\sigma_k^2/\sum_i x_i$. Every linear unbiased competitor,
$\sum_i a_i x_i=1$, has at least that variance. With at least two contributors
and nonzero individual volumes, the residual estimator is unbiased for
$\sigma_k^2$.
\end{theorem}
\begin{proof}\leanok Independence makes the linear-estimator variance a
diagonal quadratic form. Complete its square around $a_i=1/\sum_jx_j$. For the
residual estimator, subtract the fitted-mean square from the true-centre sum;
their expectations are $\sigma_k^2$ times $|A|$ and one, respectively.
\end{proof}

\begin{definition}[Practical assessments]\label{def:bfAssessments}\lean{VerifiedReserving.bfYRaw, VerifiedReserving.bfSigma2Raw, VerifiedReserving.bfYStdErrSq, VerifiedReserving.bfZStdErrSq, VerifiedReserving.bfProcessErrorEstimate, VerifiedReserving.bfEstimationErrorEstimate, VerifiedReserving.bfMsepEstimate, VerifiedReserving.bfCovarianceApprox}\leanok
Equations (3), (4), (6), and (7) define the raw development proportions, raw
column variances, squared standard errors of the selected proportions, and the
minimum of the prefix and suffix variance sums. The single-year plug-in is
\[
 \widehat{\operatorname{msep}}
 =\widehat{\operatorname{Var}}(R_i)
  +(U_i^2+s_U^2)s_z^2+s_U^2(1-z)^2.
\]
The total-reserve covariance assessment on page 99 is also a definition. Its
correlations, standard errors, smoothing, and tail values are supplied by the
actuary. No unbiasedness or consistency claim is attached to those selected
or assessed inputs.

The standalone equation (3), (4), and (6) definitions retain Mack's one-based
symbols \texttt{n} and \texttt{k}; the stochastic theorems above instead use an explicit
finite contributor set and therefore follow the repository's zero-based
triangle convention.
\end{definition}

\begin{theorem}[Assessment algebra and total covariance]\label{thm:bfTotalCovariance}\lean{VerifiedReserving.bfMsepEstimate_eq, VerifiedReserving.covariance_mul_mul_of_pair_indep}\leanok
\uses{def:bfAssessments}
The single-year assessment is exactly its named process term plus its two
estimation terms. Before Mack omits the middle term for the total-reserve
approximation, independence of the pairs $\{X,W\}$ and $\{Y,Z\}$ gives
\[
 \operatorname{Cov}(XY,WZ)
 =\operatorname{Cov}(X,W)E[Y]E[Z]
  +\operatorname{Cov}(X,W)\operatorname{Cov}(Y,Z)
  +E[X]E[W]\operatorname{Cov}(Y,Z).
\]
\end{theorem}
\begin{proof}\leanok Factor the required cross moments between the two
independent pairs, substitute the covariance identities, and normalize the
polynomial expression. \end{proof}

\begin{theorem}[A finite BF witness]\label{thm:bfWitness}\lean{VerifiedReserving.BFPredictionWitness.X, VerifiedReserving.BFPredictionWitness.bf1, VerifiedReserving.BFPredictionWitness.bf2, VerifiedReserving.BFPredictionWitness.bf3, VerifiedReserving.BFPredictionWitness.future_moments, VerifiedReserving.BFPredictionWitness.conditional_msep_zero}\leanok
\uses{def:bfModel, thm:bfModelMoments, thm:bfPredictionError}
On the one-point probability space, take $C_{i,k}=k+1$, so every increment is
one, $x_i=y_k=1$, and $\sigma_k^2=0$ on the one-period horizon. BF1--BF3 and
the normalization hold together. The future reserve has mean one and variance
zero, and prediction by one has conditional MSEP zero. This checks
satisfiability of the hypotheses; it does not provide a positive-variance BF
model.
\end{theorem}
\begin{proof}\leanok Every random variable is constant. The generated
sigma-algebras are trivial, so the full increment family is independent. \end{proof}
